\documentclass[12pt,a4paper]{article}
\usepackage[utf8]{inputenc}
\usepackage[T1]{fontenc}
\usepackage[spanish]{babel}
\usepackage[margin=1.4cm]{geometry}
\usepackage{array,tabularx}
\usepackage{enumitem}
\usepackage{titlesec}
\usepackage{xcolor}
\usepackage{hyperref}

% Restrained, high-contrast palette shared by headings, emphasis and links.
\definecolor{HeadingColor}{HTML}{243B53}
\definecolor{InstructionColor}{HTML}{704214}
\definecolor{ConceptColor}{HTML}{205C50}
\definecolor{LinkColor}{HTML}{0047AB}
\hypersetup{colorlinks=true,allcolors=LinkColor,pdfborder={0 0 0}}
\newcommand{\guidelink}[2]{\href{#1}{\textcolor{LinkColor}{\textbf{\underline{#2}}}}}
\newcommand{\instruction}[1]{\textcolor{InstructionColor}{\textbf{#1}}}
\newcommand{\concept}[1]{\textcolor{ConceptColor}{\textbf{#1}}}

\setlength{\parindent}{0pt}
\setlength{\parskip}{3pt}
\setlength{\tabcolsep}{5pt}
\renewcommand{\arraystretch}{1.12}
\setlist{leftmargin=0.6cm,topsep=3pt,partopsep=0pt,parsep=0pt,itemsep=1pt}
\titleformat{\section}{\normalsize\bfseries\color{HeadingColor}}{}{0pt}{}
\titlespacing*{\section}{0pt}{8pt}{3pt}
\newcolumntype{Y}{>{\raggedright\arraybackslash}X}
\newcommand{\checkbox}{\fbox{\rule{0pt}{0.8ex}\rule{0.8ex}{0pt}}}

\begin{document}

\begin{center}
{\Large\bfseries\color{HeadingColor} C5--C6: Organizar la información y priorizar una oportunidad de innovación.}\\[3pt]
Emprendimiento e Innovación --- Grado 10
\end{center}

\section*{Objetivo}
\concept{C5:} Organizar evidencia contextual para conectar usuarios, experiencias, espacios, necesidades y hallazgos.\\
\concept{C6:} Priorizar una oportunidad considerando relevancia, viabilidad y beneficio esperado. 

Primero comprendan el contexto; en esta etapa todavía no deben proponer soluciones.

\section*{Instrucciones de trabajo y entrega}
\instruction{Esta guía es de consulta: las respuestas van en un único documento acumulativo del proyecto en LaTeX.} Continúen el archivo del proyecto si existe; de lo contrario, créenlo en \guidelink{https://www.overleaf.com/}{Overleaf}. Usen un título estable, integrantes, grado y fecha de actualización. Conserven las fuentes de C1--C4 y añadan secciones C5 y C6; después incorporarán C7--C10 en el mismo archivo. Entreguen en cada etapa el PDF actualizado sin duplicar secciones anteriores.

La IA es opcional y solo puede apoyar el código LaTeX, el formato y la organización del contenido elaborado por ustedes. \instruction{No debe inventar evidencia, conclusiones, decisiones, ideas ni aportes de los integrantes.} Los ejemplos de esta guía son ficticios; no son resultados del grupo.

\section*{C5 --- 1. Integren el usuario y el contexto}
Retomen las observaciones y respuestas de usuarios de C1--C4. Si ya tienen un perfil o mapa, completen ese registro; no lo reconstruyan. Organicen la información en \textbf{una sola representación}: mapa, diagrama o tabla. Incluyan:
\begin{itemize}
\item La situación y el usuario principal; justifiquen su elección con un registro real.
\item Qué hace, qué necesita lograr y qué dificultades experimenta el usuario.
\item Lugar, momentos, personas y recursos relevantes, y sus relaciones.
\item Acciones antes, durante y después de la situación, señalando dónde aparece la dificultad.
\end{itemize}
No necesitan un perfil, un mapa y una tabla de recorrido separados. Elijan el formato que muestre mejor las relaciones. Si usan \guidelink{https://app.diagrams.net/}{diagrams.net}, incorporen la representación al documento LaTeX.

\section*{C5 --- 2. Construyan un registro único de hallazgos}
Seleccionen los hallazgos relevantes y asígnenles códigos H1, H2, etc. Un \concept{hallazgo} es una conclusión sustentada en los registros; no es una solución. Usen una fila por hallazgo:

\begin{tabularx}{\linewidth}{|p{0.08\linewidth}|Y|Y|Y|}
\hline
\textbf{Código} & \textbf{Hallazgo (conclusión)} & \textbf{Evidencia y fuente} & \textbf{Relación / pendiente} \\
\hline
H1, etc. & Qué aprendieron sobre usuario y contexto. & Observación o respuesta; registro, actividad y fecha disponible. & Vínculo con otro hallazgo y supuesto o pregunta pendiente, si existe. \\
\hline
\end{tabularx}

\textbf{Ejemplo ficticio:} observar una fila al inicio del descanso es evidencia; concluir que ese momento concentra la espera requiere comparar registros. Afirmar que ocurre por falta de personal es un supuesto si no se investigó esa causa.

Distingan \concept{observación o respuesta registrada}, \concept{conclusión} y \concept{supuesto por verificar}. No inventen citas ni fechas. Si falta respaldo, registren la afirmación como supuesto. Enlacen los códigos H con la representación anterior; no hagan otro mapa. Al pie del registro, señalen la pregunta pendiente más importante. No añadan una reflexión que repita esta información.

\newpage
\section*{Puente de C5 a C6}
Ya tienen el contexto organizado y hallazgos vinculados a sus fuentes. Ahora usarán esos códigos para comparar oportunidades de mejora. \instruction{C5 describe y relaciona lo investigado; C6 decide qué necesidad merece atención. Ninguna etapa selecciona soluciones.}
\section*{C6 --- 1. Delimiten las oportunidades que compararán}
Una \concept{oportunidad} es una necesidad o situación que vale la pena mejorar; todavía no es una solución. Retomen entre dos y tres oportunidades identificadas previamente. Si solo investigaron una situación, comparen dos o tres aspectos de mejora ya reconocidos en sus hallazgos, sin inventar problemas nuevos.

Describan cada opción en una frase e indiquen su usuario, necesidad y evidencia disponible. Asignen códigos (A, B, C) y registren las descripciones en las filas de la matriz de la actividad 3; debajo, relacionen cada código con sus hallazgos H y fuentes, sin crear otra tabla. Identifiquen la fuente: actividad, registro o pregunta, y fecha si está disponible. Pueden referirse a un hallazgo de C5 y al registro que lo respalda.

\textbf{Ejemplo orientativo:} ``Reducir la espera de estudiantes que compran en la cafetería durante el descanso'' expresa una oportunidad. ``Crear una aplicación de pedidos'' ya propone una solución.

\section*{C6 --- 2. Comparen con una matriz sencilla}
Usen los mismos criterios para todas las opciones. La \concept{viabilidad} se refiere a la posibilidad de abordar la oportunidad como grupo, considerando tiempo, acceso a usuarios, recursos y permisos; no exige diseñar ni presupuestar una solución.

{\small
\begin{tabularx}{\linewidth}{|p{0.17\linewidth}|Y|Y|Y|}
\hline
\textbf{Criterio} & \textbf{1: Bajo} & \textbf{2: Medio} & \textbf{3: Alto} \\
\hline
Relevancia & Dificultad ocasional y de poca importancia para los usuarios. & Dificultad recurrente o importante para algunos usuarios. & Dificultad frecuente y prioritaria para los usuarios. \\
\hline
Viabilidad & Hay barreras importantes de tiempo, acceso o permisos. & Es abordable, pero hay condiciones por confirmar. & El tiempo, acceso y recursos disponibles permiten abordarla. \\
\hline
Beneficio esperado & La mejora tendría un efecto limitado en la experiencia. & La mejora tendría un efecto apreciable en la experiencia. & La mejora tendría un efecto significativo en una necesidad importante. \\
\hline
\end{tabularx}
}

\textbf{Relevancia} pregunta cuánto importa la situación actual; \textbf{beneficio esperado}, qué cambio favorable podría producir atenderla. Este beneficio es una expectativa razonada, no un resultado demostrado ni una promesa.


\section*{C6 --- 3. Registren la comparación}
Elaboren la matriz, puntúen cada criterio de 1 a 3 y sumen: \concept{Total = relevancia + viabilidad + beneficio esperado} (máximo 9). Todos los criterios tienen el mismo peso.

\begin{tabularx}{\linewidth}{|Y|c|c|c|c|}
\hline
\textbf{Oportunidad} & \textbf{Relevancia} & \textbf{Viabilidad} & \textbf{Beneficio} & \textbf{Total} \\
\hline
A: descripción breve & 1--3 & 1--3 & 1--3 & 3--9 \\
\hline
B: descripción breve & 1--3 & 1--3 & 1--3 & 3--9 \\
\hline
\end{tabularx}

Debajo, justifiquen \textbf{cada puntaje} con evidencia o condiciones concretas. Si falta información, escriban \textbf{SD (sin datos)} y señalen qué falta: no asignen 1 por desconocimiento ni calculen un total completo para esa opción. Una opción menos investigada no es necesariamente menos importante.

\textbf{Ejemplo de cálculo, con puntajes ficticios:} A obtiene $3+2+3=8$ y B obtiene $2+3+2=7$. A tendría prioridad si sus puntajes están sustentados. No usen estos números como evidencia propia.

\newpage
\section*{C6 --- 4. Seleccionen y justifiquen una oportunidad}
Elijan \textbf{una} oportunidad y escriban un párrafo que explique:
\begin{itemize}
\item Cuál seleccionaron y cómo se compara con las otras opciones.
\item Qué observaciones o respuestas de usuarios respaldan su relevancia y beneficio esperado; incluyan al menos dos referencias concretas a la evidencia disponible.
\item Qué condiciones hacen viable abordarla y qué limitación deben reconocer.
\end{itemize}
El total orienta la decisión. Si hay empate, revisen la evidencia y las barreras de viabilidad; registren el criterio que resolvió el empate. Si eligen una opción con menor total, expliquen la razón. Si la evidencia es insuficiente, indiquen que la selección es provisional y qué dato falta, sin presentarla como una certeza.

\section*{Resultado y continuidad}
Conserven los hallazgos, sus fuentes, la matriz y la oportunidad seleccionada: serán el punto de partida de C7--C8. Remitan a los códigos H y de oportunidad. En C7 convertirán la oportunidad en un desafío abierto a distintas respuestas. No deben diseñar una solución todavía. Entreguen el PDF compilado del documento acumulativo.
\section*{Verificación final}
{\small
\begin{itemize}[label=\checkbox]
\item C5 integra usuario, contexto y recorrido en una representación y relaciona hallazgos con fuentes y pendientes.
\item Integramos mapas y tablas para evitar repetir información; distinguimos evidencia, supuestos y preguntas pendientes.
\item Comparamos dos o tres oportunidades derivadas de C5 con relevancia, viabilidad y beneficio esperado; justificamos cada puntaje.
\item Seleccionamos una oportunidad con evidencia, condiciones concretas y limitaciones reconocidas.
\item La IA no fabricó información. Compilamos y revisamos que el PDF sea legible.
\end{itemize}
}
\end{document}